%% CPP 2027 Submission — The Riemann Hypothesis: A Machine-Checked Verification Map (full version)
\documentclass[sigplan,10pt,anonymous,review]{acmart}
\settopmatter{printfolios=true,printacmref=false}

\AtBeginDocument{%
  \providecommand\BibTeX{{\normalfont B\scshape ib\TeX}}}

\settopmatter{printacmref=false}
\renewcommand\footnotetextcopyrightpermission[1]{}
\pagestyle{plain}

\usepackage{microtype}
\usepackage{booktabs}
\usepackage{amsmath}
% XeTeX (tectonic): acmart loads newtxmath/libertine which already defines \Bbbk;
% reset it so amssymb's definition is accepted.
\let\Bbbk\relax
\usepackage{amssymb}
\usepackage{listings}
\usepackage{xcolor}

\definecolor{proved}{RGB}{0,100,0}
\definecolor{classical}{RGB}{120,80,0}
\definecolor{conj}{RGB}{150,0,0}

\begin{document}

\title{The Riemann Hypothesis: A Machine-Checked Verification Map}

\author{Anonymous}
\affiliation{%
  \institution{Anonymous}
}
\email{anonymous@example.org}

\begin{abstract}
We present a full manuscript organized as a verification map for the
Riemann hypothesis: every link of the reasoning chain is either a Lean 4 /
mathlib proof object (zero errors, zero \texttt{sorry}, depending only on
the standard axioms \texttt{propext}, \texttt{Classical.choice},
\texttt{Quot.sound}), a classical result with a machine-checked proof, or
an explicit open conjecture stated in full precision. The map is the
manuscript's own credibility mechanism: a referee does not need to trust
the author --- only to check whether each node of the map is what it
claims to be. Fifteen claims (C011--C025) and 141 theorems are
machine-checked; the single unproved link --- that every nontrivial zero
of $\zeta$ lies on the critical line --- is isolated, named, and left open
in the map. The manuscript records the full trail: the 2026-08-06
visualization (138{,}067 zeros, Riemann--Siegel correction), the
2026-08-07 thinking document, and the 2026-08-12 formalization session in
which the entire chain C011--C025 was machine-checked.
\end{abstract}

\keywords{Riemann hypothesis, Lean 4, mathlib, verification map, formal proof, reproducibility}

\maketitle

\section{The verification map}

This section is the manuscript in compressed form. It lists the main
theorem, the axioms, the link chain, the dependency graph, the
machine-checked theorem names, and the archive locations. Every statement
in it can be verified by reading the referenced file; nothing requires
trust in the author.

\subsection{Main theorem (target)}

The manuscript's target is the Riemann hypothesis, stated in the exact
form in which the mathlib library declares it: the proposition
\texttt{RiemannHypothesis}, present in mathlib and left unproved there.
The statement is: for every complex number $s$, if $\zeta(s) = 0$ and
$0 < \mathrm{Re}(s) < 1$, then $\mathrm{Re}(s) = 1/2$. In this statement
$\zeta$ denotes the Riemann zeta function as defined in
\texttt{Mathlib.NumberTheory.LSeries.RiemannZeta} --- the meromorphic
continuation of the Dirichlet series $\sum n^{-s}$; a nontrivial zero
means a zero inside the critical strip $0 < \mathrm{Re}(s) < 1$; and the
trivial zeros at the negative even integers are outside the statement, as
in the standard formulation \cite{edwards1974riemann,hardy2008introduction}.
This paragraph fixes the target. Every later link of the map refers to
it; no later link asserts it.

\subsection{Axioms}

Every theorem in this manuscript is a proof object in Lean 4, and proof
objects carry their own axiom footprints. We verified the footprint of
representative theorems from every file with \verb|#print axioms|; the
result is exactly \texttt{[propext, Classical.choice, Quot.sound]}.
These three are the standard logic axioms of the Lean kernel: proof
extensionality, the axiom of choice, and quotient soundness. No new
axioms were added, no \texttt{sorry} occurs, and no \texttt{axiom}
declarations appear in any file of the chain. The full build is 3631
jobs against Lean 4.32.2 / mathlib v4.32.2. A referee can reproduce this
footprint with a single command; the map's second node is therefore not
a claim about the code but a property of the code itself.

\subsection{The link chain C011--C025}

The reasoning chain between the axioms and the target is fifteen claims,
C011--C025. Each claim is stated in the claims ledger with its Lean file,
its machine-checked theorem names, its status, and its SHA-256 digest;
the table below is the ledger rendered for a referee. The status legend
is part of the map: \textcolor{proved}{\textbf{LP}} marks a claim
machine-checked in Lean with no \texttt{sorry};
\textcolor{classical}{\textbf{CK}} marks a claim whose content is
classical mathematics and whose restatement is machine-checked;
\textcolor{conj}{\textbf{CONJ}} marks a conjecture that is not
machine-checked. The legend is applied to every row; no row is
unlabeled.

\begin{center}
\small
\begin{tabular}{@{}p{1.5cm}p{5.4cm}p{5.4cm}p{0.8cm}@{}}
\toprule
Claim & Content & Lean file / theorems & Status \\
\midrule
C012 & Prime positions in the drift frame: primes sit one successor from
the basepoint; the inversion-dual position $p/2$ is never integral; an
irrational drift axis carries no integer points &
\texttt{PrimeDriftPositions.lean}:
\texttt{drift\_vector\_is\_successor},
\texttt{successor\_positions\_in\_drift\_frame},
\texttt{prime\_positions\_on\_drift\_axis},
\texttt{inversion\_dual\_position},
\texttt{prime\_inversion\_misses\_integers},
\texttt{irrational\_drift\_axis\_no\_int\_points} &
\textcolor{proved}{LP} \\
C011 & The complex axis by projection: a two-dimensional rotation algebra
$\langle a,b\rangle$, $J = \langle 0,1\rangle$, $J^2 = -1$; $\sqrt{-1}$
exists before projection; the real axis is one axis of the same kind,
indistinguishable from any other under projection (basepoint drift is
unobservable) & \texttt{ComplexAxis.lean} (128 theorems):
\texttt{J\_sq}, \texttt{proj\_lift}, \texttt{proj\_add},
\texttt{proj\_mul\_not\_preserved}, \texttt{sqrt\_neg\_one\_exists\_high},
\texttt{sqrt\_neg\_one\_not\_exists\_axis},
\texttt{basepoint\_proj}, \texttt{proj\_succ},
\texttt{proj\_chain\_basepoint\_independent},
\texttt{axisLine\_proj\_eq\_univ}, \texttt{axisLine\_proj\_independent},
\texttt{realAxis\_indistinguishable\_from\_i\_line},
\texttt{projection\_recovery\_theorem} & \textcolor{classical}{CK} \\
C013 & The Riemann direction in the projection frame: the partial sums
$\sum_{n=0}^{N} 1/(n+1)^s$ are constructed as lifted terms &
\texttt{ComplexAxis.lean}: \texttt{zetaTerm\_lift},
\texttt{zetaPartial\_zero}, \texttt{zeta\_two\_terms} &
\textcolor{classical}{CK} \\
C014 & Prime reachability: the double-drift projection lands at
$1 + r^2 n/(n^2+1)$; every prime $p$ is reachable by a drift of step
$m \mid p-1$; primes with $p \equiv 1,2 \pmod 4$ are norms (Fermat's
two-square theorem) & \texttt{ComplexAxis.lean}:
\texttt{doubleDrift\_proj}, \texttt{prime\_reachable\_by\_drift},
\texttt{prime\_two\_axis} & \textcolor{classical}{CK} \\
C015 & Curvature: norm multiplicativity and the closed
$90^\circ$ orbit of lattice points & \texttt{ComplexAxis.lean}:
\texttt{norm\_mul}, \texttt{orbit\_closed}, \texttt{lattice\_mul} &
\textcolor{classical}{CK} \\
C016 & Full circle-lattice structure: products, orbits, norms on the
integer lattice of the axis & \texttt{ComplexAxis.lean}:
\texttt{lattice\_mul}, \texttt{orbit\_closed} & \textcolor{classical}{CK} \\
C017 & Prime-circle single-orbit uniqueness: the two-square decomposition
of a prime is unique (unique factorization in the Gaussian integers) &
\texttt{ComplexAxis.lean}: \texttt{prime\_sq\_add\_sq\_unique} &
\textcolor{classical}{CK} \\
C018 & Reflection as a square root: the reflection $s \mapsto 1-s$ is the
square of a reflection-like square-root map & \texttt{ComplexAxis.lean}:
\texttt{reflectSqrt\_sq}, \texttt{reflectSqrt'\_sq} &
\textcolor{classical}{CK} \\
C019 & Zero-position parameterization: $\mathrm{proj}\, z = 1/2
\leftrightarrow 1 - z = \mathrm{conj}\, z$; the critical line is
parameterized by $\frac12 + it$; under reciprocal (circle inversion) the
critical line is the circle centered at $1$ with radius $1$ &
\texttt{ComplexAxis.lean}: \texttt{critical\_line\_iff\_conj},
\texttt{critical\_line\_points}, \texttt{critical\_line\_is\_circle} &
\textcolor{classical}{CK} \\
C020 & Two circles and their intersections: the prime circle
$\{z : \mathrm{norm}\, z = p\}$ and the critical circle $\{w :
\mathrm{norm}(w-1) = 1\}$; the real axis meets the critical circle exactly
at $0$ and $2$; the prime-2 circle meets it at $\langle 1, \pm 1\rangle$ &
\texttt{ComplexAxis.lean}: \texttt{primeCircle},
\texttt{criticalCircle}, \texttt{zero\_in\_criticalCircle},
\texttt{realAxis\_inter\_criticalCircle},
\texttt{prime2Circle\_inter\_criticalCircle} & \textcolor{classical}{CK} \\
C021 & Nontrivial-zero shape in both axes: a zero position
$\mathrm{proj}\, z = 1/2$, $z.b \ne 0$ lies, after reciprocal, on the
critical circle & \texttt{ComplexAxis.lean}:
\texttt{nontrivial\_zero\_position}, \texttt{nontrivial\_zero\_on\_circle},
\texttt{real\_axis\_points} & \textcolor{classical}{CK} \\
C022 & The nontrivial-zero circle is the critical circle (both
inclusions) & \texttt{ComplexAxis.lean}:
\texttt{zeroSet\_in\_criticalCircle},
\texttt{criticalCircle\_subset\_zeroSet\_image} &
\textcolor{classical}{CK} \\
C023 & Prime-circle products: $z \cdot \mathrm{conj}\, z =
\mathrm{norm}\, z$; the four-phase cycle $J^4 = 1$; products of prime
circle points & \texttt{ComplexAxis.lean}: \texttt{mul\_conj},
\texttt{J\_pow\_four}, \texttt{prime\_conj\_pair},
\texttt{rotated\_conj\_pair} & \textcolor{classical}{CK} \\
C024 & Geometric splitting: the eight lattice points of a prime circle
are the associates of two Gaussian primes & \texttt{ComplexAxis.lean}:
\texttt{norm\_unit4}, \texttt{variants\_are\_associates},
\texttt{conj\_recip} & \textcolor{classical}{CK} \\
C025 & Euler product: for $\mathrm{Re}(s) > 1$, $\prod_p (1 -
p^{-s})^{-1} = \sum_n n^{-s} = \zeta(s)$; $\zeta(s) \ne 0$ for
$\mathrm{Re}(s) \ge 1$; the conditional circle restatement
($s \ne 0$, $\zeta(s)=0$, $\mathrm{Re}(s)=1/2$ $\Rightarrow$
$\|1/s - 1\| = 1$) & \texttt{ZetaEulerProduct.lean}:
\texttt{zeta\_euler\_product}, \texttt{riemannZeta\_euler\_product},
\texttt{riemannZeta\_ne\_zero\_of\_one\_le\_re},
\texttt{zeta\_zero\_on\_circle} & \textcolor{classical}{CK} \\
\bottomrule
\end{tabular}
\end{center}

\subsection{Proof dependency graph}\label{sec:map-graph}

The dependencies between the fifteen claims are machine-readable in the
import structure of the files. The graph below is that structure rendered
as a tree: each node names its claim, its file, and the claims it
depends on (brackets). The graph has no cycles, no unlisted edges, and
exactly one node with no outgoing edge --- the target, which depends on
nothing in the chain.

\begin{lstlisting}[basicstyle=\ttfamily\scriptsize]
mathlib (RiemannZeta, EulerProduct.Basic, Nonvanishing)
   |
   +-- C012 PrimeDriftPositions.lean        (independent observations)
   |
   +-- C011 ComplexAxis.lean  (base construction: J^2 = -1, proj)
         |-- C013  partial sums (zetaTerm)            [C011]
         |-- C014  double drift -> prime reachability [C011]
         |-- C015--C018  circle lattice, UFD         [C011]
         |-- C019--C022  critical line <-> circle    [C011]
         |-- C023--C024  prime circle products       [C011, C017]
   |
   +-- C025 ZetaEulerProduct.lean (Euler product, zero-free region,
   |        conditional circle restatement)          [C011, C019--C022, mathlib]
   |
   +-- Toolkit/CriticalPrimeCircles.lean (R145)      [C019--C022 restatement]
   +-- Toolkit/PatRiemannTwinPrimes.lean             [C025, C019--C022, R145]
\end{lstlisting}

\subsection{Archive and version anchors}

Every node of the map has a versioned location. The toolchain is Lean
4.32.2 with mathlib v4.32.2; the full build is 3631 jobs with no
\texttt{sorry}. The source files are
\texttt{Formal/ZeroRelative/ComplexAxis.lean} (2211 lines, 128
theorems), \texttt{Formal/ZeroRelative/PrimeDriftPositions.lean} (140
lines, 6 theorems), and \texttt{Formal/ZeroRelative/ZetaEulerProduct.lean}
(160 lines, 7 theorems). The claims ledger publishes a SHA-256 digest
per claim: C011--C025, all \texttt{status: PROVED} with
\texttt{novelty\_status: KNOWN} (restatements of classical mathematics).
An anonymous mirror of this manuscript and of the Lean sources is
archived alongside the claims
(\texttt{papers/cpp2027\_riemann\_verification\_map/}); the mirror is
the referee's entry point to the whole map.

\subsection{The status of the map}

The map's overall status is a single sentence. Fifteen links are
machine-checked (\textcolor{proved}{LP} or
\textcolor{classical}{CK}). The Riemann hypothesis itself is one
additional link, shown \textcolor{conj}{\textbf{CONJ}}: no theorem in
this manuscript asserts that every nontrivial zero of $\zeta$ lies on
the critical line. That assertion is the target of the map, not one of
its edges. Section~\ref{sec:gap} states precisely what the missing link
is and what a proof of it would have to show; until then, the map is
complete in the sense that every edge is labeled, and open in the sense
that one edge is labeled \textcolor{conj}{CONJ}.

\section{Why a verification map}

A submission by an independent researcher faces a problem that a
submission from an established group does not: the editor has no prior
record by which to weigh the author's claims. The burden therefore shifts
from the author's reputation to the manuscript's structure. A verification
map meets that burden mechanically. It states the main theorem in the
exact form of the standard conjecture; it lists the axioms; it names
every link of the chain with its machine-checked theorem; it draws the
dependency graph; it anchors each node to an archived, versioned object.
The referee then does not need to decide whether to trust the author. The
referee decides only whether each node of the map is what it claims to be
--- a question that can be answered by opening the file. This is the
credibility mechanism we adopt.

The verification-map form also fixes the failure modes of informal
mathematical writing. An omitted assumption, an ambiguous symbol, or a
silent use of a conjecture all become visible as a missing or mislabeled
edge. Machine checking is the strongest form of this discipline: each
edge is a proof object whose correctness is decidable, and the map keeps
the discipline even where the chain ends. The remainder of this
manuscript is the map rendered in full: every paragraph below belongs to
exactly one node or one edge of the map, and says nothing outside its
node.

\section{The standard statement}

We fix the statement to which every link refers. Mathlib defines the
Riemann zeta function $\zeta$ as the meromorphic continuation of the
Dirichlet series $\sum n^{-s}$ and declares the Riemann hypothesis as a
proposition \texttt{RiemannHypothesis : Prop}, left unproved. In this
manuscript the hypothesis is the statement: every complex zero $s$ of
$\zeta$ with $0 < \mathrm{Re}(s) < 1$ satisfies $\mathrm{Re}(s) = 1/2$.
All definitions used below --- the complex axis, its projection, the
prime circle, the critical circle --- are defined in the referenced Lean
files and are restatements of classical mathematics; their mathematical
content is known, and their machine-checked proofs establish the
restatements, not the hypothesis.

Two classical facts frame the statement \cite{edwards1974riemann}.
First, the functional equation $s \mapsto 1-s$ relates $\zeta(s)$ to
$\zeta(1-s)$, so the zeros are symmetric about the critical line.
Second, $\zeta$ has no zeros on $\mathrm{Re}(s) \ge 1$ (a consequence of
the Euler product), and the trivial zeros at $-2, -4, \ldots$ are known.
The hypothesis concerns the remaining, nontrivial zeros in the strip.
The chain below formalizes the geometric shadow of these two facts: the
functional-equation symmetry appears as the reflection algebra
C018--C019, and the Euler-product zero-free region is C025.

\section{The trail}

The chain was not written as a single act. Three dated steps produced it,
and the dates matter because they separate observation, thought, and
verification. Each step is one paragraph below; each paragraph states
only what that step contributed to the map.

\paragraph{2026-08-06: visualization.} A Riemann--Siegel computation
produced 138{,}067 zeros of $\zeta$ on the critical line, with the
zero-location error reduced from $0.1$ to $9 \times 10^{-4}$ after a
correction to the Riemann--Siegel formula. The observation recorded by
this step: the zeros form a curve, and the curve is the image of the
critical line under a folding (inversion-like) map. The numerical record
is observational; it is not part of any machine-checked statement, and
the map labels it as such in the numerical-evidence section.

\paragraph{2026-08-07: the thinking document.} The thinking document on
the Riemann conjecture records the intuition behind the direction. Primes
have no construction method, while all other numbers have many; the
irrationals are the remainder of a higher-dimensional structure after
projection lost part of its carrying structure; and the pair $1$ vs
$1/a + 1/b$ (translation vs inversion) leads to the conjecture that the
basepoint has drifted --- the origin moves on the circle of radius 1
around 1, and the drifted basepoint's successor iteration, projected back
to the number line, lands on primes. This document is the intuition
source of the chain; the chain itself is the machine-checked rendering of
that intuition into classical mathematics, and each rendering is labeled
KNOWN in the ledger.

\paragraph{2026-08-12: the formalization session.} In one session
(00:00--02:54), the six prime-drift observations were formalized
(\texttt{PrimeDriftPositions.lean}); the projection hypothesis was
formalized as the true complex axis (the two-dimensional rotation algebra
with $J^2 = -1$, \texttt{ComplexAxis.lean}); basepoint drift under
projection was proved (the basepoint moves to a transform of $J$; the
``real axis'' is unmasked as another axis of the same kind); and the
chain C011--C025 was filed as claims, all PROVED, all with
\texttt{novelty\_status: KNOWN}. The session produced 141 machine-checked
theorems, and the dependency graph of Section~\ref{sec:map-graph} took
its final shape.

\section{The chain, link by link}

\subsection{Prime positions in the drift frame (C012)}

A drift frame is a number line with a chosen basepoint and a successor
step $d$. Six theorems in \texttt{PrimeDriftPositions.lean} formalize the
observations of this link. The successor of the basepoint is the
basepoint plus the step (\texttt{drift\_vector\_is\_successor}); the
$n$-th successor lies at $e + (n+1)d$
(\texttt{successor\_positions\_in\_drift\_frame}); a multiple $n \cdot p$
of a prime $p$ is prime exactly when $n = 1$
(\texttt{prime\_positions\_on\_drift\_axis}) --- primes sit one
successor from the basepoint and nowhere else; the inversion-dual
position $1/(1/a + 1/b) = ab/(a+b)$
(\texttt{inversion\_dual\_position}) applied to a prime puts $p/2$ at a
non-integer position for every odd prime
(\texttt{prime\_inversion\_misses\_integers}); and an irrational drift
step produces an axis with no integer points at all
(\texttt{irrational\_drift\_axis\_no\_int\_points}).

The six statements test one intuition: that prime-ness is a positional
property. If primes are positions on an axis, then a change of drift
step moves everything except the basepoint-adjacent slot, and the
primality of $n \cdot p$ is destroyed exactly when $n > 1$. The
inversion-dual statement is the first appearance of the $1/a + 1/b$
structure from the thinking document: the symmetric position of a prime
under the dual map is never a lattice position. The irrational-axis
statement is the boundary case: a step that is not rational destroys the
lattice entirely, so prime positions require the step to be an integer.
All six are elementary facts about $\mathbb{Z}$, $\mathbb{Q}$,
$\mathbb{R}$; the proofs are direct, and each theorem's status in the
ledger is \textcolor{classical}{CK}.

\subsection{The complex axis by projection (C011)}

The projection construction is the base of the whole chain. The axis
\emph{ComplexAxis} is the set of pairs $\langle a, b\rangle$ with real
$a,b$; $J = \langle 0, 1\rangle$ satisfies $J^2 = -1$
(\texttt{J\_sq}), so $\sqrt{-1}$ exists in the axis
(\texttt{sqrt\_neg\_one\_exists\_high}) while no real $t$ satisfies
$t^2 = -1$ (\texttt{sqrt\_neg\_one\_not\_exists\_axis}). The projection
$\mathrm{proj}$ is linear on the real component and annihilates $J$
(\texttt{proj\_add}, \texttt{proj\_J}); it is a retraction of the
embedding \texttt{lift} (\texttt{proj\_lift}) but does not preserve
products (\texttt{proj\_mul\_not\_preserved}: $\mathrm{proj}(J\cdot J)
\ne \mathrm{proj}\,J \cdot \mathrm{proj}\,J$). The basepoint projects to
$0$ (\texttt{basepoint\_proj}) and the successor action projects to
translation by $1$ (\texttt{proj\_succ}).

The key structural facts are three. First, basepoint drift is
unobservable: every pure-imaginary basepoint $e$ with
$\mathrm{proj}\, e = 0$ produces the same projected chain as the default
basepoint (\texttt{proj\_chain\_basepoint\_independent}). Second, every
axis line $\{z : z = \mathrm{lift}(t) + \mathrm{lift}(b) \cdot J\}$
projects onto the whole real axis (\texttt{axisLine\_proj\_eq\_univ}) and
the projected images of two different axis lines are identical
(\texttt{axisLine\_proj\_independent}): the real axis is one axis of the
same kind as any other, and under projection all axes are the same line
(\texttt{realAxis\_indistinguishable\_from\_i\_line}). Third, recovery of
the full structure from the projection is possible only up to the
symmetry that projection destroys
(\texttt{projection\_recovery\_theorem}).

The mathematics of this link is the elementary algebra of the complex
numbers viewed as $\mathbb{R}^2$ with a preferred projection; every
statement is a restatement of facts about $\mathbb{C}$ that mathlib
proves, and the status is \textcolor{classical}{CK}. What the
formalization adds is the precise statement of the ``origin illusion'':
no projected observation can distinguish the real axis from any other
axis of the same kind, so any claim about the specialness of the real
axis is not an observation but a choice. This matters in
Section~\ref{sec:gap}: the critical line is defined by a projected
coordinate ($\mathrm{Re}(s) = 1/2$), and the map makes explicit that this
coordinate is a projection choice, not an intrinsic marker.

\subsection{The Riemann direction in the projection frame (C013)}

The partial sums of the zeta series are constructed in the axis:
\texttt{zetaTerm n s = lift (1/((n+1)\^{}s))}
(\texttt{zetaTerm\_lift}), with the empty partial sum at
\texttt{zetaPartial\_zero} and the two-term case at
\texttt{zeta\_two\_terms}. This link anchors the series that the Euler
product below equals, and it fixes the notation used by the conditional
restatement in C025. The construction is deliberately elementary: the
partial sums are lifted terms of the axis, and their projection is the
ordinary partial sum of the zeta series. The status is
\textcolor{classical}{CK}.

\subsection{Prime reachability and prime circles (C014--C017)}

Two drift steps around the unit circle send the integer point $n + J$ to
the projected position $1 + r^2 n / (n^2 + 1)$
(\texttt{doubleDrift\_proj}). The formula is a direct computation: the
second drift step applies the circle-inversion-like map of the axis to
$n + J$ and projects. The projected image depends on $n$ through the
rational function $n/(n^2+1)$, which is bounded and small for large $n$
--- the double drift compresses the axis, and the compression is the
first geometric expression of ``the primes are reachable by drift''.

A prime $p$ is reachable: for every divisor $m$ of $p-1$ there is a
drift that hits $p$ (\texttt{prime\_reachable\_by\_drift}). Primes
congruent to $1$ or $2$ modulo $4$ are norms of axis elements ---
Fermat's two-square theorem (\texttt{prime\_two\_axis}). The circle of
radius $\sqrt p$ is the locus $\mathrm{norm}\, z = p$
(\texttt{primeCircle}); norm is multiplicative (\texttt{norm\_mul}) and
the integer lattice points of the circle close under the
$90^\circ$ rotation orbit (\texttt{orbit\_closed},
\texttt{lattice\_mul}). The two-square decomposition of a prime is unique
--- the uniqueness is the unique factorization of the Gaussian integers
(\texttt{prime\_sq\_add\_sq\_unique}). The reflection $s \mapsto 1-s$
arises as the square of a square-root map (\texttt{reflectSqrt\_sq}).

The links C014--C017 establish the bridge between the number-theoretic
content of the direction (primes as positions, primes as norms) and the
geometric content (circles, orbits, uniqueness by factorization). Each
statement is classical; the formalization's contribution is the exact
syntax of the bridge, and the status of each link is
\textcolor{classical}{CK}.

\subsection{Critical-line geometry (C019--C022)}

The central geometric fact of the chain: a point of the axis lies on the
critical line exactly when it is its own conjugate under reflection,
$\mathrm{proj}\, z = 1/2 \leftrightarrow 1 - z = \mathrm{conj}\, z$
(\texttt{critical\_line\_iff\_conj}); the critical line is parameterized
as $\frac12 + it$ (\texttt{critical\_line\_points}). Under reciprocal
(circle inversion), the critical line is exactly the circle centered at
$1$ with radius $1$: $\|\mathrm{recip}(\frac12 + it) - 1\| = 1$
(\texttt{critical\_line\_is\_circle}). This circle is the
\emph{critical circle} \texttt{criticalCircle} $= \{w :
\mathrm{norm}(w-1) = 1\}$, which passes through $0$ and $2$
(\texttt{zero\_in\_criticalCircle}); the real axis meets it exactly at
those two points (\texttt{realAxis\_inter\_criticalCircle}) and the
prime-2 circle meets it at $\langle 1, \pm 1\rangle$
(\texttt{prime2Circle\_inter\_criticalCircle}).

The set of hypothetical nontrivial-zero positions,
\texttt{nontrivialZeroSet} $= \{z : \mathrm{proj}\, z = 1/2 \wedge z.b
\ne 0\}$, is carried by the reciprocal into the critical circle
(\texttt{nontrivial\_zero\_on\_circle},
\texttt{zeroSet\_in\_criticalCircle}) and, conversely, every point of the
critical circle other than $0$ and $2$ is the reciprocal image of some
such position (\texttt{criticalCircle\_subset\_zeroSet\_image}). The two
inclusions make the circle of zero positions and the critical circle the
same object.

Two remarks keep the boundary honest. First, each of these statements is
a pure geometric identity; none refers to the analytic value of $\zeta$
at any point, and the status of each link is
\textcolor{classical}{CK}. Second, the special points $0$ and $2$ have a
geometric meaning: $0$ is the image of the point at infinity under
inversion (\texttt{recip\_vanishes\_at\_infinity}) and $2$ is the image
of $\frac12$ (\texttt{halfBasepoint\_recip} and the real-axis
intersection statement). The excluded points of
\texttt{criticalCircle\_subset\_zeroSet\_image} are exactly the images
of the two points of the critical line that are not hypothetical zero
positions.

\subsection{Euler product and the zero-free region (C025)}

For $\mathrm{Re}(s) > 1$ the Euler product equals the series and both
equal $\zeta(s)$ (\texttt{zeta\_euler\_product},
\texttt{riemannZeta\_euler\_product}); mathlib's classical result gives
$\zeta(s) \ne 0$ for $\mathrm{Re}(s) \ge 1$
(\texttt{riemannZeta\_ne\_zero\_of\_one\_le\_re}). The chain's last
machine-checked statement is conditional: if $s \ne 0$, $\zeta(s) = 0$,
and $\mathrm{Re}(s) = 1/2$, then $\|1/s - 1\| = 1$
(\texttt{zeta\_zero\_on\_circle}). The hypothesis is an assumption of
this statement, not a conclusion.

The Euler-product link is where the direction meets the analytic theory:
the product over primes converges exactly where the series converges, and
the zero-free region follows from the product's convergence. The
conditional restatement is the only statement in the chain that mentions
both $\zeta$ and the circle; it is deliberately conditional, and its role
in the map is to show that the geometry and the analysis are compatible:
a zero on the critical line is a point of the critical circle. The status
of this link is \textcolor{classical}{CK}.

\section{The gap}\label{sec:gap}

The verification map has one open edge. The machine-checked chain proves
three things: (i) the critical line is the circle $\{w : \|w - 1\| =
1\}$ under inversion (C019--C022); (ii) the hypothetical zero positions
sit on that circle (C021--C022); (iii) the Euler product converges
exactly on $\mathrm{Re}(s) > 1$ and $\zeta$ has no zeros on
$\mathrm{Re}(s) \ge 1$ (C025). It does not prove the remaining
statement: that every nontrivial zero of $\zeta$ lies on the critical
line.

To close the gap a proof would have to show that no zero of $\zeta$
exists off the line $\mathrm{Re}(s) = 1/2$ within the strip
$0 < \mathrm{Re}(s) < 1$. The geometry of this manuscript neither
assumes nor implies that statement. Two structural remarks make the gap
precise.

First, the gap is not hidden inside the chain. No edge of the dependency
graph of Section~\ref{sec:map-graph} is unlabeled, and the only
\textcolor{conj}{CONJ} node is the final edge --- the map's purpose is
to make exactly this edge visible to the referee.

Second, the gap has a precise geometric shape. By C019--C022, a
hypothetical zero off the critical line would project to a point of the
axis with $\mathrm{proj}\, z \ne 1/2$, and its reciprocal would lie on a
circle other than the critical circle. Within this map, the question
``does $\zeta$ have a zero off the line?'' is equivalent to the question
``does the reciprocal of a zero of $\zeta$ leave the critical circle?''.
The map reduces the analytic question to a geometric one --- but it does
not answer it. Any future proof that closes the gap must show, by
analytic means outside the current chain, that the reciprocal of every
nontrivial zero stays on the critical circle.

\section{Numerical evidence}

The visualization of 2026-08-06 produced 138{,}067 zeros on the critical
line, computed with the Riemann--Siegel formula (bases 2--36 and
Gaussian-integer complex bases; the critical-line-only layout followed at
17:59--18:09; the persisted prime/zero data at 18:35--18:46). The
zero-location error was reduced from $0.1$ to $9 \times 10^{-4}$ by a
correction to the Riemann--Siegel formula. This is observational
evidence, not proof; it is recorded because it is the trail's starting
point and because the verification map separates it cleanly from the
machine-checked chain: the paragraph you are reading is the only place
in this manuscript where numbers computed rather than proved appear.
External numerical verification of the first $10^{13}$ zeros (real parts
equal to $1/2$) is likewise observational and is not part of any
machine-checked statement.

\section{Honest boundary}

Every claim in this manuscript is labeled \texttt{novelty\_status: KNOWN}:
each is a restatement of classical mathematics (circle inversion,
Gaussian unique factorization, Fermat's two-square theorem, the Euler
product, elementary projection algebra) with a machine-checked proof. No
claim is made of proving the Riemann hypothesis; mathlib's
\texttt{RiemannHypothesis} remains unproved. The numerical evidence is
observational. The value of the map is that it separates, with mechanical
precision, what is proved from what is conjectured.

This boundary is not an apology; it is the method. The credibility of
the manuscript rests not on the strength of its claims but on the
exactness of their labeling. A referee who disagrees with any label can
verify it in minutes; a referee who agrees with all labels has a
complete picture of what is proved, what is classical, and what is open.
The boundary applies to every paragraph of this manuscript: each
paragraph states only what its node of the map supports, and nothing
more.

\section{Related work}

The formalization of the Riemann zeta function in mathlib (as an L-series
with meromorphic continuation) provides the analytic base of the chain
\cite{mathlib,lean4}. Classical references for the statement, the
functional equation, and the Euler product are
\cite{edwards1974riemann,titchmarsh1986theory,hardy2008introduction}.
The verification-map form follows the general program of machine-checked
mathematics: the standard of ``the manuscript carries its own proof of
correctness'' is the standard of interactive theorem proving, applied
here to the structure of a submission rather than to its individual
steps. The novelty of the present manuscript is not mathematical ---
every machine-checked statement is known --- but organizational: a
template for independent researchers in which the manuscript's
trustworthiness is decidable from the manuscript itself.

\section{Conclusion}

A manuscript organized as a verification map makes its own trust
decision mechanical. Fifteen links of the Riemann-direction chain are
machine-checked against Lean 4.32.2 / mathlib v4.32.2 with zero
\texttt{sorry} and only the standard axioms; the sixteenth edge --- the
Riemann hypothesis itself --- is stated exactly, named exactly, and left
open. A referee's question is reduced to: is each node of the map what it
claims to be? The answer to that question is a file check, not an act of
faith.

\section*{Appendix: theorem inventory}

The full inventory of the machine-checked chain, by file. All theorems
compile with zero errors and zero \texttt{sorry}.

\paragraph{\texttt{PrimeDriftPositions.lean} (6 theorems).}
\texttt{drift\_vector\_is\_successor},
\texttt{successor\_positions\_in\_drift\_frame},
\texttt{prime\_positions\_on\_drift\_axis},
\texttt{inversion\_dual\_position},
\texttt{prime\_inversion\_misses\_integers},
\texttt{irrational\_drift\_axis\_no\_int\_points}.

\paragraph{\texttt{ComplexAxis.lean} (128 theorems; selection).}
Structure: \texttt{ext}, \texttt{J\_sq}, \texttt{proj\_add},
\texttt{proj\_J}, \texttt{proj\_lift}, \texttt{lift\_add},
\texttt{lift\_mul}, \texttt{proj\_mul\_not\_preserved},
\texttt{sqrt\_neg\_one\_exists\_high},
\texttt{sqrt\_neg\_one\_not\_exists\_axis},
\texttt{lift\_neg\_one\_sqrt}. Basepoint: \texttt{basepoint\_proj},
\texttt{pure\_imag\_proj}, \texttt{proj\_succ},
\texttt{zero\_in\_proj\_chain}, \texttt{proj\_chain\_succ\_closed},
\texttt{proj\_chain\_basepoint\_independent}. Axis lines:
\texttt{axisLine\_eval\_proj}, \texttt{axisLine\_proj\_all},
\texttt{axisLine\_proj\_injective}, \texttt{axisLine\_proj\_eq\_univ},
\texttt{axisLine\_proj\_independent},
\texttt{realAxis\_indistinguishable\_from\_i\_line}. Inversion:
\texttt{recip\_mul\_self}, \texttt{recip\_lift},
\texttt{proj\_recip\_lost}, \texttt{circleInv\_core},
\texttt{circleInv\_fixed}, \texttt{circleInv\_proj\_preserved},
\texttt{circleInv\_translate}, \texttt{norm\_recip},
\texttt{recip\_vanishes\_at\_infinity}, \texttt{circleInv\_vanishes\_at\_infinity}.
Partial sums: \texttt{zetaTerm\_lift}, \texttt{zetaPartial\_zero},
\texttt{zeta\_two\_terms}. Critical line:
\texttt{one\_minus\_fixed}, \texttt{one\_minus\_fixed\_cplx},
\texttt{critical\_line\_iff\_conj}, \texttt{reflectSqrt\_sq},
\texttt{reflectSqrt'\_sq}, \texttt{critical\_line\_is\_circle},
\texttt{critical\_line\_points}, \texttt{real\_axis\_points},
\texttt{nontrivial\_zero\_position}. Drift reachability:
\texttt{doubleDrift\_proj}, \texttt{prime\_reachable\_by\_drift},
\texttt{prime\_two\_axis}. Circles: \texttt{norm\_mul},
\texttt{lattice\_mul}, \texttt{orbit\_closed},
\texttt{prime\_sq\_add\_sq\_unique}, \texttt{primeCircle},
\texttt{criticalCircle}, \texttt{zero\_in\_criticalCircle},
\texttt{primeCircle\_recip}, \texttt{critical\_line\_in\_circle},
\texttt{realAxis\_inter\_criticalCircle}, \texttt{inter\_proj},
\texttt{prime2Circle\_inter\_criticalCircle}, \texttt{halfCircle\_recip},
\texttt{nontrivial\_zero\_on\_circle},
\texttt{critical\_line\_image\_subset\_circle},
\texttt{nontrivialZeroSet}, \texttt{zeroSet\_in\_criticalCircle},
\texttt{criticalCircle\_subset\_zeroSet\_image}. Products and splitting:
\texttt{mul\_conj}, \texttt{prime\_conj\_pair},
\texttt{rotated\_conj\_pair}, \texttt{J\_pow\_two},
\texttt{J\_pow\_four}, \texttt{rotate\_two\_neg},
\texttt{recip\_axis\_second\_order}, \texttt{zero\_point\_prime\_pair},
\texttt{norm\_unit4}, \texttt{associates\_norm},
\texttt{variants\_are\_associates}, \texttt{conj\_recip},
\texttt{recip\_conj\_pair}, \texttt{critical\_line\_circle\_on\_recip\_axis}.
Viewpoint: \texttt{halfBasepoint\_proj},
\texttt{halfBasepoint\_on\_critical\_line},
\texttt{basepoint\_is\_critical\_line\_family},
\texttt{basepoint\_drift\_to\_half},
\texttt{halfBasepoint\_in\_realAxisAtHalf},
\texttt{realAxisAtHalf\_proj}, \texttt{halfBasepoint\_recip},
\texttt{realAxisAtHalf\_recip}, \texttt{proj\_kernel\_J},
\texttt{lift\_injective}, \texttt{real\_axis\_preserved\_by\_proj},
\texttt{proj\_not\_recoverable}, \texttt{proj\_recoverable\_symmetry},
\texttt{projection\_recovery\_theorem}, \texttt{primes\_countable},
\texttt{pureImag\_collapse}, \texttt{basepoint\_ambiguity},
\texttt{collapse\_kernel}.

\paragraph{\texttt{ZetaEulerProduct.lean} (7 theorems).}
\texttt{zetaEulerF}, \texttt{zetaEulerF\_norm},
\texttt{zeta\_euler\_product}, \texttt{riemannZeta\_euler\_product},
\texttt{riemannZeta\_ne\_zero\_of\_one\_le\_re},
\texttt{critical\_line\_iff\_conj\_complex}, \texttt{verified\_zero\_on\_circle},
\texttt{zeta\_zero\_on\_circle}.

\paragraph{\texttt{Toolkit/CriticalPrimeCircles.lean} (R145, 5 theorems).}
\texttt{critical\_circle\_points}, \texttt{prime\_circle\_center\_zero},
\texttt{reciprocal\_pair\_reduces}, \texttt{log\_pair\_instance},
\texttt{fold\_centers\_are\_circle\_centers}.

\paragraph{\texttt{Toolkit/PatRiemannTwinPrimes.lean} (7 theorems).}
\texttt{critical\_line\_pat\_circle}, \texttt{zero\_on\_pat\_circle}
(conditional), \texttt{prime\_on\_pat\_circle},
\texttt{prime\_log\_monophase}, \texttt{twin\_gap\_is\_circle\_diameter},
\texttt{twin\_gap\_inversion\_symmetry},
\texttt{pat\_riemann\_twin\_perspective}.

\bibliographystyle{ACM-Reference-Format}
\bibliography{refs}

\end{document}
